\documentclass[10pt]{article}
\usepackage[hypertexnames=false,colorlinks=true,breaklinks]{hyperref}
\usepackage{url}
\usepackage{geometry}
\usepackage{enumitem}

\geometry{top=0.6in,bottom=0.6in,left=1.75in,right=0.6in}
\pagestyle{empty}

\setlength{\parindent}{0mm}
\setlength{\parskip}{4pt}

% Custom command to recreate your marginpar style compactly without pushing lines down
% \newcommand{\sectionleft}[1]{\noindent\makebox[0pt][r]{\makebox[1.5in][l]{\textbf{#1}}\hspace{0.25in}}}
\newcommand{\sectionleft}[1]{\noindent\makebox[0pt][r]{\makebox[1.1in][l]{\textbf{#1}}\hspace{0.2in}}}

\begin{document}

{\huge \textbf{Joaqu\'{i}n Rapela}}

\vspace{0.05in}
\textit{\textbf{Address}}: Gatsby Computational Neuroscience Unit, University College London, 25 Howland St, London W1T 4JG, UK
\textit{\textbf{Phone}}: +44 7432049398 \quad \textit{\textbf{Email}}: \href{mailto:j.rapela@ucl.ac.uk}{j.rapela@ucl.ac.uk} \quad \textit{\textbf{WWW}}: \url{http://www.gatsby.ucl.ac.uk/~rapela}
\textit{\textbf{GitHub}}: \url{https://github.com/joacorapela} \quad \textit{\textbf{Right to work}}: UK Indefinite Leave to Remain (ILR) -- No sponsorship required
\textit{\textbf{Last update}}: \today

\vspace{0.15in}

\sectionleft{Education}%
\textbf{University of Southern California} \hfill 2003---2010 \\
PhD in Electrical Engineering. \\
Advisors: Dr. Norberto M. Grzywacz (Neuroscience) and Dr. Jerry M. Mendel (Signal Processing) \\
Neuroscience Graduate Program (Studies in Neuroscience) \hfill 2001---2003 \\
MS in Electrical Engineering \hfill 2000---2003

\vspace{0.05in}
\sectionleft{}%
\textbf{Universidad de Buenos Aires} \hfill 1992---1998 \\
``Licenciatura (equivalent to MS)'' and BS in Computer Science.

\vspace{0.15in}

\sectionleft{Publications}%
Campagner D., \ldots, \textbf{Rapela J.}, \ldots, SWC GCNU Experimental
Neuroethology Group (2025). \textit{Aeon: an open-source platform to study the
neural basis of ethological behaviours over naturalistic timescales}.
\href{https://www.biorxiv.org/content/10.1101/2025.07.31.664513v1}{bioRxiv
2025.07.31.664513}.

\vspace{0.05in}
\textbf{Rapela J.}, Westerfield M., Townsend J. (2018). \textit{A new
foreperiod effect on single-trial phase coherence. Part I: existence and
relevance}.
\href{https://www.mitpressjournals.org/doi/abs/10.1162/neco_a_01109}{Neural
Computation 30(9):2348-2383}.

\vspace{0.05in}
\sectionleft{}%
\textbf{Rapela J.} (2018). \textit{Traveling waves appear and disappear in
unison with produced speech}. \href{https://arxiv.org/abs/1806.09559}{arXiv:1806.09559 [q-bio.NC]}.

\vspace{0.05in}
\sectionleft{}%
\textbf{Rapela J.}, Kostuk M, Rowat, P.F., Mullen, T., Chang E.F., Bouchard K.
(2015) \textit{Modeling neural activity at the ensemble level}.
\href{http://arxiv.org/abs/1505.00041}{arXiv:1505.00041 [q-bio.NC]}.

\vspace{0.05in}
\sectionleft{}%
\textbf{Rapela J.}, Felsen G., Touryan J., Mendel J.M., \& Grzywacz N.M. (2010). \textit{ePPR: a new strategy for the characterization of sensory cells from input/output data}.
\href{https://www.tandfonline.com/doi/full/10.3109/0954898X.2010.488714}{Network:
Computation in Neural Systems. 21(1-2): 35-90}.

\vspace{0.05in}
\sectionleft{}%
\textbf{Rapela J.}, Mendel J.M., \& Grzywacz N.M. (2006). \textit{Estimating
nonlinear receptive fields from natural images}.
\href{https://jov.arvojournals.org/article.aspx?articleid=2192869}{Journal of
Vision 6(4), 441-474}.

\vspace{0.05in}
\sectionleft{}%
Shattuck D., \textbf{Rapela, J.}, Asma E., Chatziioannou A., Qi J., \& Leahy R. (2002).
\textit{Internet2-based 3D PET Image Reconstruction using a PC Cluster}.
\href{https://iopscience.iop.org/article/10.1088/0031-9155/47/15/317}{Physics
in Medicine and Biology 47, 2785-2795}.

\vspace{0.05in}
\sectionleft{}%
Full publication list available on \href{https://scholar.google.com/citations?user=eXkDg2UAAAAJ&hl=en}{Google Scholar}.

\vspace{0.15in}

\sectionleft{Academic}%
\textbf{Gatsby Computational Neuroscience Unit, University College London} \hfill June 2019---Present \\
\sectionleft{Experience}%
\textit{Position: Research Engineer Fellow}
\begin{itemize}[leftmargin=1.5em, topsep=2pt, itemsep=1pt, parsep=1pt]
    \item Develops, optimises, and maintains open-source implementations of state-of-the-art machine learning algorithms for probabilistic inference in high-dimensional neural and behavioural time series (e.g., \href{https://github.com/joacorapela/svGPFA}{svGPFA}, \href{https://github.com/joacorapela/ssm}{Linear Dynamical Systems}) utilizing Python, JAX, PyTorch, and Plotly.
    \item Spearheads curriculum development and delivers lectures on statistical neuroscience and neuroinformatics for PhD candidates at the Sainsbury Wellcome Centre, UCL (\href{https://github.com/joacorapela/statNeuro2025}{Statistical Neuroscience 2025}, \href{https://github.com/joacorapela/neuroinformatics24}{Neuroinformatics 2024}).
    \item Secured grant funding, directed the hiring process, and managed a team of research software engineers to execute the BBSRC-funded project \href{https://gtr.ukri.org/projects?ref=BB\%2FW019132\%2F1}{Machine Intelligence for Neuroscience Experimental Control}.
    \item Mentored undergraduate fellows in statistical neuroscience for the \href{https://www.simonsfoundation.org/grant/shenoy-undergraduate-research-fellowship-in-neuroscience-surfin/}{SURFiN Program}: Aisha Qureshi (QMUL, CS, 2021--2022); Zimo Li (UCL, NPP, 2023--2024).
    \item Collaborates with experimental neuroscientists at the SWC on advanced statistical methods for characterising neural populations electrophysiological recordings.
    \item Co-develops the data analytics pipeline for the \href{https://www.biorxiv.org/content/10.1101/2025.07.31.664513v1}{Aeon platform}, enabling storage, management, visualization, and high-throughput statistical analysis of continuous (24/7) behavioural and electrophysiological mouse datasets.
    \item Designs and delivers training workshops on open science and software best practices for researchers at Gatsby and SWC (e.g., introduction to Python, code testing CI pipelines).
    \item Collaborates with members of SWC and Gatsby to enhance institutional research culture through the \href{https://sainsburywellcomecentre.github.io/RCWG/}{Research Culture Working Group}.
\end{itemize}

\vspace{0.08in}

\sectionleft{}%
\textbf{The Truccolo Lab for Computational Neuroscience, Brown University} \hfill Aug 2017---Jan 2019 \\
\sectionleft{}%
\textit{Position: Research Associate}
\begin{itemize}[leftmargin=1.5em, topsep=2pt, itemsep=1pt, parsep=1pt]
    \item Discovered stereotypical discrete states in pre-ictal, ictal, and post-ictal periods of seizures separated by more than one hour, utilizing micro-electrode array layer 2/3 cortical recordings from patients with epilepsy.
    \item Extracted two-dimensional descriptors of high-dimensional electrophysiological recordings from patients with epilepsy that robustly separate functional ictal periods using t-SNE, XGBoost, and manifold learning.
\end{itemize}

% \clearpage % Control page breaks gracefully 

\sectionleft{}%
\textbf{Swartz Center for Computational Neuroscience, UC San Diego} \hfill Nov 2010---July 2017 \\
\sectionleft{}%
\textit{Position: Postdoctoral and Visiting Scholar}
\begin{itemize}[leftmargin=1.5em, topsep=2pt, itemsep=1pt, parsep=1pt]
    \item Characterized spatio-temporal brain dynamics related to speech production from electrocorticographic (ECoG) neural recordings, analyzing synchronization and phase-amplitude coupling.
    \item Served as Principal Investigator for a project funded by the Center for Brain Activity Mapping (part of the B.R.A.I.N. initiative), evaluating computations of neural ensembles using high-resolution recordings.
    \item Discovered a new effect of temporal expectation on single-trial phase coherence of EEG data using a novel variational Bayes linear regression method.
    \item Found that switching attention between visual and auditory modalities generates transient arousal of attention across both modalities via time-frequency analysis.
    \item Developed a computer game controlled online by eye movements (electrooculography) to quantify the speed and accuracy of attention-orienting mechanics.
\end{itemize}

\vspace{0.08in}

\sectionleft{}%
\textbf{University of Southern California} \hfill 2001---2010 \\
\sectionleft{}%
\textit{Position: Research Assistant (Advisor: Dr. Norberto M. Grzywacz \& Dr. Jerry Mendel)}
\begin{itemize}[leftmargin=1.5em, topsep=2pt, itemsep=1pt, parsep=1pt]
    \item Investigated Bayesian models to evaluate the mathematical optimality of retinal codes.
    \item Developed the Volterra Relevant Space Technique (VRST) for structural spatial Volterra model estimation.
    \item Developed the extended Projection Pursuit Regression (ePPR) algorithm for structural characterization of visual cells.
\end{itemize}

\vspace{0.05in}
\sectionleft{}%
\textbf{University of Southern California} \hfill Aug 2000---Aug 2001 \\
\sectionleft{}%
\textit{Position: Research Assistant (Advisor: Dr. Richard M. Leahy)} \vspace{-4pt}
\begin{itemize}[leftmargin=1.5em, topsep=2pt, itemsep=1pt, parsep=1pt]
    \item Co-developed a distributed computing Java-based system for 3D PET data reconstruction and performed source localization.
\end{itemize}

\vspace{0.15in}

\sectionleft{Teaching}%
\textbf{University College London} (Course Organiser and Instructor) \hfill 2019---2025 \\
\sectionleft{Experience}%
Instructor for \textit{PyStarters} (2019, 2021).\\
Teaching Assistant for \href{https://github.com/joacorapela/neuroinformatics23}{NEURO019: Neuroinformatics} (2023).\\
Module organizer and instructor for the \href{https://github.com/joacorapela/gcnuBridging2023}{Probability Module of the Gatsby Bridging Program} (2023). \\
Course organiser and instructor for \href{https://github.com/joacorapela/neuroinformatics24}{Neuroinformatics for SWC PhD Students} (2024). \\
Course organiser and instructor for \href{https://github.com/joacorapela/statNeuro2025}{Statistical Neuroscience for SWC PhD Students} (2025).

\vspace{0.05in}
\sectionleft{}%
\textbf{Brown University} \hfill 2017 \\
Teaching Assistant for \textit{NEUR2110: Statistical Neuroscience}.

\vspace{0.05in}
\sectionleft{}%
\textbf{University of Southern California} \hfill 2009 \\
Discussion Section Leader for \textit{EE 364: Introduction to Probability and Statistics for EE/CS}.

\vspace{0.05in}
\sectionleft{}%
\textbf{Universidad de Buenos Aires} \hfill 1998---1999 \\
Teaching Assistant for \textit{Numerical Linear Algebra} and \textit{Artificial Intelligence}.

\vspace{0.15in}

\sectionleft{Service}%
\textbf{Journal Ad-hoc Reviewer:} \textit{Vision Research}, \textit{Frontiers in Neuroscience}, \textit{Journal of Perceptual Imaging}. \\
\textbf{Grant Reviewer:} The Software Sustainability Institute (SSI).

\vspace{0.15in}

\sectionleft{Funding}%
\textbf{Biotechnology and Biological Sciences Research Council (BBSRC)} \hfill 2023---2026 \\
Co-Investigator, Grant \texttt{BB/W019132/1}, \textbf{£800,000}. \\
Project: \textit{Machine Intelligence for Neuroscience Experimental Control.}

\vspace{0.05in}
\sectionleft{}%
\textbf{Center for Brain Activity Mapping (CBAM)} \hfill 2013---2014 \\
Principal Investigator, Grant \texttt{2013-023CBAM}, \textbf{\$30,000}. \\
Project: \textit{Taking the next step toward understanding computations by neural ensembles with high resolution neural recordings.}

% \clearpage

\sectionleft{Selected}%
\textbf{Mathematics:} Topology, Fundamental Concepts of Analysis, Real Analysis, Functional Analysis. \\
\sectionleft{Courses}%
\textbf{Engineering:} Transform Theory, Random Processes, Digital Signal Processing, Optimization, Information Theory. \\
\textbf{Neuroscience:} Control and Communication in the Nervous System, Neurobiology, Advanced Neurosciences I \& II.

\vspace{0.15in}

\sectionleft{Industry}%
\textbf{IBM Almaden Research Center} (Staff Software Engineer) \hfill Jan 2000---Aug 2000 \\
\sectionleft{Experience}%
\textbf{Oracle, Argentina} (Software Engineer) \hfill June 1998---Sept 1998 \\
\textbf{IBM, Argentina \& IBM Almaden} (Research Fellow) \hfill Oct 1996---June 1998 \\
\textbf{IBM, Argentina} (Fellowship Holder) \hfill April 1994---Oct 1996

\vspace{0.15in}

\sectionleft{Programming}%
Python, R, C\#, Matlab, Java, C/C++, Smalltalk, Pascal, Assembler.
% \sectionleft{Languages}%

\end{document}
